\section{An explicit size test for independent cells}

The original model and its definition of head count are unchanged.
Let a fixed cell consist of $M$ Boolean inputs in dimension $n$.
Suppose that every independent assignment of signs on that cell is the
restriction of a function computed by $H$ original heads.  Parameters may
be chosen after seeing all labels.  Assume $n,H\geq1$ and $H\leq M$.
Then the kernel-checked theorem
\texttt{cellShatteredByHeads\_log\_size} gives
\[
 M\leq30nH\log(8eM).
\]
The proof restricts the original cleared score polynomials to the cell.
Their degree is at most $H+1$ in $1+2H(n+1)$ real parameters, and every
strict sign pattern occurs.  The repository's proved Warren bound gives
the displayed inequality; no new axiom or alternative head class is used.
The parameter-larger-than-sample case follows directly and is included.

Consequently, if $nHG\leq M$ and $G\geq60\log(8e)$, then
\[
 \exp(G/60)\leq M.
\]
This is \texttt{cellShatteredByHeads\_exponential\_size}.
A gain $G=n^\delta$ therefore requires an exponentially large local
domain in $n^\delta$.  A gain $G=\log(n)g(n)$ requires
$M\geq n^{g(n)/60}$ in the same explicit range.  In particular, a
uniform polynomial bound on the local domain cannot produce an
unbounded improvement beyond the logarithmic saving by independent
compilation alone.

For a finite family of cells with sizes $M_j\leq L$, suppose their union
covers the cube, and every cell independently admits all labels using
$H_j$ heads with $1\leq H_j\leq M_j$.  Even with overlaps,
\[
 2^n\leq\sum_jM_j
 \leq30n\log(8eL)\sum_jH_j.
\]
The theorem \texttt{independentCellFamily\_budget} checks the second
inequality and takes the first as its explicit cover-count hypothesis.
Overlaps can only increase the sum of cell sizes; they are not silently
counted as new inputs.

These statements are necessary conditions for the stated independent
cell scheme.  They do not lower-bound a construction whose heads jointly
serve different cells, whose encoding changes with labels, or whose
subproblems have correlated target labels.  The present run continues to
investigate precisely those possibilities.

All four declarations are in
\texttt{formalization/HeadComplexity/Results/IndependentCellSize.lean}.
Their transitive axiom output contains exactly \texttt{propext},
\texttt{Classical.choice}, and \texttt{Quot.sound}.
